\chapter{Zeitplan}

\begin{table}[h]
	\centering
	\begin{tabular}{|c|p{10cm}|c|}
		\hline
		\textbf{Woche} & \textbf{Aufgaben} & \textbf{Dauer} \\
		\hline
		1 & \textbf{Themenfindung und Planung} & 1 Woche \\
		& - Detaillierte Planung der Arbeit und Erstellung eines groben Konzepts & \\
		\hline
		2-3 & \textbf{Recherche und theoretische Grundlagen} & 2 Wochen \\
		& - Literaturrecherche zu Teststrategien und Datenbanken & \\
		& - Grundlagen zu Spring Boot und React/TypeScript & \\
		\hline
		4-6 & \textbf{Entwicklung der Backend-Module (Spring Boot)} & 3 Wochen \\
		& - Einrichtung des Spring Boot Projekts (Initializr) & \\
		& - Entwicklung von Kern-APIs (z. B. Benutzerregistrierung, Vertragsmanagement) & \\
		& - Anbindung an H2 für schnelle Tests & \\
		\hline
		7-8 & \textbf{Frontend-Entwicklung (React/TypeScript)} & 2 Wochen \\
		& - Einrichtung des Projekts mit Create React App (CRA) & \\
		& - Implementierung der Benutzeroberfläche für Kern-Features & \\
		\hline
		9-10 & \textbf{Integration Frontend und Backend} & 2 Wochen \\
		& - API-Anbindung und Kommunikation testen & \\
		& - Fehlerbehebung und Optimierung & \\
		\hline
		11 & \textbf{Testdurchführung und Datensammlung} & 1 Woche \\
		& - Entwicklung und Ausführung von Testszenarien & \\
		& - Sammlung von Daten zur Laufzeit und Fehlererkennung & \\
		\hline
		12 & \textbf{Datenanalyse, Visualisierung und Schreiben} & 1 Woche \\
		& - Auswertung der Ergebnisse, Diagramme erstellen & \\
		& - Fertigstellung des Berichts und Überarbeitung & \\
		\hline
		\textbf{Puffer} & \textbf{Zusätzliche Woche für Integration und Tests} & 1 Woche \\
		\hline
	\end{tabular}
	\caption{Grobzeitplan für die Bachelorarbeit}
\end{table}
